Bodový graf závislosti Giniho koeficientu příjmové nerovnosti na pokrytí \ac{KS} potvrzuje zápornou korelaci: vyšší \ac{KS} pokrytí je asociováno s~nižší příjmovou nerovností --- vztah, který je v~empirické literatuře o~institucích trhu práce považován za jeden z~nejrobustnějších.
Mechanismus je přímočarý: odvětvové \ac{KS} s~tarify závaznými pro celý sektor (erga omnes) stlačují spodní část mzdové distribuce, zvyšují mezd zaměstnanců ve spodním kvartilem, a~tím snižují roztažení distribuce měřené Ginim.
ČR s~pokrytím cca \SI{35}{\percent} a~aktuálním Gini cca 27 se nachází v~prostoru grafu, kde by hypotetické zvýšení pokrytí na \SI{70}{\percent} (středoevropský cíl) odpovídalo poklesu Gini o~2--3 body --- změna, která by CZ přiblížila finskému nebo belgickému vzoru.
Redistributivní efekt \ac{KV} institucí je tedy empiricky prokazatelný i~bez zvyšování přímých daní nebo transferů: jde o~tržní redistribuci prostřednictvím vyjednávání, která je z~hlediska politické ekonomie snáze prosaditelná než fiskální redistribuce.

Bodový graf závislosti Giniho koeficientu příjmové nerovnosti na pokrytí \ac{KS} potvrzuje zápornou korelaci: vyšší \ac{KS} pokrytí je asociováno s~nižší příjmovou nerovností --- vztah, který je v~empirické literatuře o~institucích trhu práce považován za jeden z~nejrobustnějších.
Mechanismus je přímočarý: odvětvové \ac{KS} s~tarify závaznými pro celý sektor (erga omnes) stlačují spodní část mzdové distribuce, zvyšují mezd zaměstnanců ve spodním kvartilem, a~tím snižují roztažení distribuce měřené Ginim.
ČR s~pokrytím cca \SI{35}{\percent} a~aktuálním Gini cca 27 se nachází v~prostoru grafu, kde by hypotetické zvýšení pokrytí na \SI{70}{\percent} (středoevropský cíl) odpovídalo poklesu Gini o~2--3 body --- změna, která by CZ přiblížila finskému nebo belgickému vzoru.
Redistributivní efekt \ac{KV} institucí je tedy empiricky prokazatelný i~bez zvyšování přímých daní nebo transferů: jde o~tržní redistribuci prostřednictvím vyjednávání, která je z~hlediska politické ekonomie snáze prosaditelná než fiskální redistribuce.

Bodový graf závislosti Giniho koeficientu příjmové nerovnosti na pokrytí \ac{KS} potvrzuje zápornou korelaci: vyšší \ac{KS} pokrytí je asociováno s~nižší příjmovou nerovností --- vztah, který je v~empirické literatuře o~institucích trhu práce považován za jeden z~nejrobustnějších.
Mechanismus je přímočarý: odvětvové \ac{KS} s~tarify závaznými pro celý sektor (erga omnes) stlačují spodní část mzdové distribuce, zvyšují mezd zaměstnanců ve spodním kvartilem, a~tím snižují roztažení distribuce měřené Ginim.
ČR s~pokrytím cca \SI{35}{\percent} a~aktuálním Gini cca 27 se nachází v~prostoru grafu, kde by hypotetické zvýšení pokrytí na \SI{70}{\percent} (středoevropský cíl) odpovídalo poklesu Gini o~2--3 body --- změna, která by CZ přiblížila finskému nebo belgickému vzoru.
Redistributivní efekt \ac{KV} institucí je tedy empiricky prokazatelný i~bez zvyšování přímých daní nebo transferů: jde o~tržní redistribuci prostřednictvím vyjednávání, která je z~hlediska politické ekonomie snáze prosaditelná než fiskální redistribuce.

Bodový graf závislosti Giniho koeficientu příjmové nerovnosti na pokrytí \ac{KS} potvrzuje zápornou korelaci: vyšší \ac{KS} pokrytí je asociováno s~nižší příjmovou nerovností --- vztah, který je v~empirické literatuře o~institucích trhu práce považován za jeden z~nejrobustnějších.
Mechanismus je přímočarý: odvětvové \ac{KS} s~tarify závaznými pro celý sektor (erga omnes) stlačují spodní část mzdové distribuce, zvyšují mezd zaměstnanců ve spodním kvartilem, a~tím snižují roztažení distribuce měřené Ginim.
ČR s~pokrytím cca \SI{35}{\percent} a~aktuálním Gini cca 27 se nachází v~prostoru grafu, kde by hypotetické zvýšení pokrytí na \SI{70}{\percent} (středoevropský cíl) odpovídalo poklesu Gini o~2--3 body --- změna, která by CZ přiblížila finskému nebo belgickému vzoru.
Redistributivní efekt \ac{KV} institucí je tedy empiricky prokazatelný i~bez zvyšování přímých daní nebo transferů: jde o~tržní redistribuci prostřednictvím vyjednávání, která je z~hlediska politické ekonomie snáze prosaditelná než fiskální redistribuce.

\input{texparts/python/coverage_gini_scatter}



